\documentclass[11pt,letterpaper]{article}

\usepackage[margin=1in]{geometry}
\usepackage{amsmath,amssymb,amsthm}
\usepackage{booktabs}
\usepackage{graphicx}
\usepackage{hyperref}
\usepackage{cleveref}
\usepackage{algorithm}
\usepackage{algpseudocode}
\usepackage{xcolor}
\usepackage{enumitem}
\usepackage{caption}
\usepackage{subcaption}
\usepackage{listings}
\usepackage{microtype}

\hypersetup{
  colorlinks=true,
  linkcolor=blue!70!black,
  citecolor=green!50!black,
  urlcolor=blue!60!black,
}

\lstset{
  language=Python,
  basicstyle=\ttfamily\small,
  keywordstyle=\color{blue!70!black}\bfseries,
  commentstyle=\color{green!50!black}\itshape,
  stringstyle=\color{red!60!black},
  numbers=left,
  numberstyle=\tiny\color{gray},
  breaklines=true,
  frame=single,
  framesep=3pt,
  rulecolor=\color{gray!40},
}

\newcommand{\ncsim}{\texttt{ncsim}}
\newcommand{\dBm}{\,\text{dBm}}
\newcommand{\dB}{\,\text{dB}}
\newcommand{\mW}{\,\text{mW}}
\newcommand{\MHz}{\,\text{MHz}}
\newcommand{\GHz}{\,\text{GHz}}
\newcommand{\Mbps}{\,\text{Mbps}}
\newcommand{\MBps}{\,\text{MB/s}}

\title{802.11 CSMA/CA Wireless Link and Interference Model\\for the ncsim Networked Compute Simulator}
\author{Bhaskar Krishnamachari\\
Autonomous Networks Research Group (ANRG)\\
University of Southern California}
\date{February 2026}

\begin{document}
\maketitle

\begin{abstract}
This document describes the 802.11 CSMA/CA wireless link and interference
model implemented in the \ncsim{} discrete-event simulator.  The model
encompasses three layers: (1)~a physical layer based on log-distance path loss
with MCS rate adaptation from IEEE~802.11n/ac/ax tables, (2)~a MAC layer
incorporating Bianchi's saturation throughput analysis for CSMA/CA contention,
and (3)~a spatial reuse model using conflict graphs derived from carrier
sensing range under the protocol interference model.  Two model variants are
provided: a static clique-based fair-share model (\texttt{csma\_clique}) and a
dynamic SINR-aware model with Bianchi MAC efficiency
(\texttt{csma\_bianchi}).  All equations presented here correspond exactly to
the software implementation.
\end{abstract}

\tableofcontents
\newpage

%% ============================================================
\section{Introduction and Motivation}
\label{sec:intro}

The \ncsim{} simulator models networked computing scenarios where computational
tasks are distributed across nodes and connected by communication links.  In
wireless deployments---particularly those using IEEE~802.11 (WiFi)---the
effective throughput of each link depends on:

\begin{enumerate}[itemsep=2pt]
\item \textbf{Distance-dependent SNR}: Path loss determines the received
  signal quality, which selects the modulation and coding scheme (MCS) and
  thus the PHY data rate.
\item \textbf{CSMA/CA contention}: Links sharing the same spectrum coordinate
  via carrier sensing; nearby links cannot transmit simultaneously, reducing
  per-link throughput.
\item \textbf{Hidden terminal interference}: Links outside each other's
  carrier sensing range may transmit simultaneously, causing SINR degradation
  at the receiver.
\end{enumerate}

Prior to this model, \ncsim{} used a simple proximity-based interference model
(factor $= 1/k$ for $k$ nearby active links) with statically specified link
bandwidths.  The model described here replaces both with physically-grounded
computations derived from RF parameters and IEEE~802.11 specifications.

%% ============================================================
\section{RF Configuration Parameters}
\label{sec:rfconfig}

The model is parameterized by an immutable \texttt{RFConfig} structure with the
following fields:

\begin{table}[h]
\centering
\caption{RF configuration parameters and defaults.}
\label{tab:rfconfig}
\begin{tabular}{@{}llll@{}}
\toprule
\textbf{Parameter} & \textbf{Symbol} & \textbf{Default} & \textbf{Description} \\
\midrule
\texttt{tx\_power\_dBm} & $P_\text{tx}$ & 20 & Transmit power (dBm) \\
\texttt{freq\_ghz} & $f$ & 5.0 & Carrier frequency (GHz) \\
\texttt{path\_loss\_exponent} & $n$ & 3.0 & Path loss exponent \\
\texttt{noise\_floor\_dBm} & $N_0$ & $-95$ & Noise floor (dBm) \\
\texttt{cca\_threshold\_dBm} & $\theta_\text{CCA}$ & $-82$ & CCA detect threshold (dBm) \\
\texttt{channel\_width\_mhz} & $W$ & 20 & Channel width (MHz) \\
\texttt{wifi\_standard} & --- & \texttt{ax} & MCS table: \texttt{n}, \texttt{ac}, \texttt{ax} \\
\texttt{shadow\_fading\_sigma} & $\sigma_\text{SF}$ & 0.0 & Shadow fading std.\ dev.\ (dB) \\
\texttt{rts\_cts} & --- & \texttt{false} & RTS/CTS enabled \\
\bottomrule
\end{tabular}
\end{table}

These parameters can be specified in the scenario YAML file under
\texttt{config.rf} and overridden via CLI flags (\texttt{--tx-power},
\texttt{--freq}, \texttt{--path-loss-exponent}, \texttt{--wifi-standard},
\texttt{--rts-cts}).

%% ============================================================
\section{Physical Layer Model}
\label{sec:phy}

\subsection{Log-Distance Path Loss}
\label{sec:pathloss}

The path loss model uses the standard log-distance formulation with a
Friis free-space reference at distance $d_0 = 1$\,m:

\begin{equation}
\label{eq:friis}
\text{PL}(d_0) = 20 \log_{10}\!\left(\frac{4\pi d_0 f}{c}\right)
\end{equation}

where $f$ is the carrier frequency in Hz and $c = 3 \times 10^8$\,m/s.  At
5\,GHz with $d_0 = 1$\,m, this gives PL$(d_0) \approx 46.4$\,dB.  At
2.4\,GHz, PL$(d_0) \approx 40.0$\,dB.

For distances $d > d_0$, the path loss grows logarithmically:

\begin{equation}
\label{eq:pathloss}
\text{PL}(d) = \text{PL}(d_0) + 10\,n \log_{10}\!\left(\frac{d}{d_0}\right)
\end{equation}

where $n$ is the path loss exponent ($n = 2$ for free space, $n = 3$--$4$ for
typical indoor/outdoor environments).

\textbf{Implementation notes:}
\begin{itemize}[itemsep=2pt]
\item For $d \leq 0$: PL $= 0$ (collocated nodes, no loss).
\item For $0 < d < d_0$: $d$ is clamped to $d_0$ (near-field effects are not modeled).
\end{itemize}

\subsection{Received Power}
\label{sec:rxpower}

The received power at distance $d$ is:

\begin{equation}
\label{eq:rxpower}
P_\text{rx}(d) = P_\text{tx} - \text{PL}(d) - X_\text{SF}
\end{equation}

where $X_\text{SF}$ is the shadow fading component in dB (see
\Cref{sec:shadow}).  All quantities are in dBm.

\subsection{Signal-to-Noise Ratio (SNR)}
\label{sec:snr}

The SNR at the receiver, in the absence of co-channel interference, is:

\begin{equation}
\label{eq:snr}
\text{SNR} = P_\text{rx} - N_0 \qquad (\text{in dB})
\end{equation}

This determines the baseline PHY rate for each link.

\subsection{Signal-to-Interference-plus-Noise Ratio (SINR)}
\label{sec:sinr}

When hidden terminal interferers are present, SINR replaces SNR.
Interference powers must be summed in the \emph{linear} (milliwatt) domain:

\begin{equation}
\label{eq:sinr}
\text{SINR} = 10 \log_{10}\!\left(
  \frac{P_\text{rx}^\text{(linear)}}{N_0^\text{(linear)} + \sum_{j \in \mathcal{H}} P_j^\text{(linear)}}
\right)
\end{equation}

where $\mathcal{H}$ is the set of active hidden terminal links (those
\emph{not} in the conflict graph neighborhood---see \Cref{sec:conflict}), and
the conversion between dBm and milliwatts is:

\begin{equation}
P^\text{(linear)} = 10^{P^\text{(dBm)}/10}
\end{equation}

\textbf{Key insight:} Only hidden terminals contribute to SINR degradation.
Links within the CSMA contention domain do \emph{not} transmit simultaneously
(carrier sensing prevents it), so they affect throughput via time-sharing, not
via interference power.  This separation is central to the
\texttt{csma\_bianchi} model (see \Cref{sec:bianchi_model}).

\subsection{Shadow Fading}
\label{sec:shadow}

When $\sigma_\text{SF} > 0$, log-normal shadow fading is applied.  For each
ordered node pair $(i, j)$, a fading value is drawn from:

\begin{equation}
X_{ij} \sim \mathcal{N}(0, \sigma_\text{SF}^2)
\end{equation}

The fading map is:
\begin{itemize}[itemsep=2pt]
\item \textbf{Symmetric:} $X_{ij} = X_{ji}$ (reciprocal channel).
\item \textbf{Deterministic:} Generated from a seeded PRNG (\texttt{random.Random(seed)}),
  ensuring identical fading across runs with the same seed.
\item \textbf{Per-node-pair:} Each unique pair $(i, j)$ with $i < j$ gets one
  independent draw.
\end{itemize}

When $\sigma_\text{SF} = 0$ (the default), the shadow fading map is empty and
$X_\text{SF} = 0$ for all links.

%% ============================================================
\section{MCS Rate Adaptation}
\label{sec:mcs}

The PHY data rate is selected from an MCS (Modulation and Coding Scheme) table
based on the link SNR.  The implementation includes tables for three 802.11
standards, each for 1 spatial stream at 20\,MHz channel width with long guard
interval:

\subsection{MCS Tables}

\begin{table}[h]
\centering
\caption{802.11ax (HE) MCS table, 20\,MHz, 1 spatial stream, 3.2\,$\mu$s GI.}
\label{tab:mcs_ax}
\begin{tabular}{@{}cllr@{}}
\toprule
\textbf{MCS} & \textbf{Modulation} & \textbf{Min SNR (dB)} & \textbf{Rate (Mbps)} \\
\midrule
0  & BPSK 1/2      & 5  & 8.6 \\
1  & QPSK 1/2      & 8  & 17.2 \\
2  & QPSK 3/4      & 11 & 25.8 \\
3  & 16-QAM 1/2    & 14 & 34.4 \\
4  & 16-QAM 3/4    & 18 & 51.6 \\
5  & 64-QAM 2/3    & 22 & 68.8 \\
6  & 64-QAM 3/4    & 25 & 77.4 \\
7  & 64-QAM 5/6    & 29 & 86.0 \\
8  & 256-QAM 3/4   & 32 & 103.2 \\
9  & 256-QAM 5/6   & 35 & 114.7 \\
10 & 1024-QAM 3/4  & 38 & 129.0 \\
11 & 1024-QAM 5/6  & 41 & 143.4 \\
\bottomrule
\end{tabular}
\end{table}

\begin{table}[h]
\centering
\caption{802.11ac (VHT) MCS table, 20\,MHz, 1 SS.}
\label{tab:mcs_ac}
\begin{tabular}{@{}cllr@{}}
\toprule
\textbf{MCS} & \textbf{Modulation} & \textbf{Min SNR (dB)} & \textbf{Rate (Mbps)} \\
\midrule
0  & BPSK 1/2     & 5  & 6.5 \\
1  & QPSK 1/2     & 8  & 13.0 \\
2  & QPSK 3/4     & 11 & 19.5 \\
3  & 16-QAM 1/2   & 14 & 26.0 \\
4  & 16-QAM 3/4   & 18 & 39.0 \\
5  & 64-QAM 2/3   & 22 & 52.0 \\
6  & 64-QAM 3/4   & 25 & 58.5 \\
7  & 64-QAM 5/6   & 29 & 65.0 \\
8  & 256-QAM 3/4  & 32 & 78.0 \\
9  & 256-QAM 5/6  & 35 & 86.7 \\
\bottomrule
\end{tabular}
\end{table}

\begin{table}[h]
\centering
\caption{802.11n (HT) MCS table, 20\,MHz, 1 SS, 800\,ns GI.}
\label{tab:mcs_n}
\begin{tabular}{@{}cllr@{}}
\toprule
\textbf{MCS} & \textbf{Modulation} & \textbf{Min SNR (dB)} & \textbf{Rate (Mbps)} \\
\midrule
0 & BPSK 1/2    & 5  & 6.5 \\
1 & QPSK 1/2    & 8  & 13.0 \\
2 & QPSK 3/4    & 11 & 19.5 \\
3 & 16-QAM 1/2  & 14 & 26.0 \\
4 & 16-QAM 3/4  & 18 & 39.0 \\
5 & 64-QAM 2/3  & 22 & 52.0 \\
6 & 64-QAM 3/4  & 25 & 58.5 \\
7 & 64-QAM 5/6  & 29 & 65.0 \\
\bottomrule
\end{tabular}
\end{table}

\subsection{Rate Selection Algorithm}

Given a link SNR value, the highest MCS index whose minimum SNR requirement is
met is selected:

\begin{equation}
\label{eq:mcs_select}
R_\text{PHY} = \max\{r_k : \text{SNR} \geq \text{SNR}_{\min,k}\}
\end{equation}

where $(r_k, \text{SNR}_{\min,k})$ are the rate and threshold for MCS index
$k$.  If the SNR is below the minimum MCS threshold (MCS~0), the rate is 0
(link not viable).

For channel widths other than 20\,MHz, rates scale linearly:

\begin{equation}
\label{eq:cw_scale}
R_\text{PHY}(W) = R_\text{PHY}(20) \times \frac{W}{20}
\end{equation}

This is approximate but standard practice for capacity estimation.

\subsection{Unit Conversion}

The simulator uses MB/s (megabytes per second) internally.  PHY rates are
converted from Mbps:

\begin{equation}
\label{eq:unit}
R\;[\MBps] = \frac{R\;[\Mbps]}{8}
\end{equation}

%% ============================================================
\section{Per-Link PHY Rate Computation}
\label{sec:phyrate}

For each link $\ell$ connecting transmitter node $i$ to receiver node $j$,
the PHY rate is computed as follows:

\begin{align}
d_{ij}    &= \sqrt{(x_i - x_j)^2 + (y_i - y_j)^2} \label{eq:dist} \\
P_\text{rx} &= P_\text{tx} - \text{PL}(d_{ij}) - X_{ij} \label{eq:prx_link} \\
\text{SNR}_\ell &= P_\text{rx} - N_0 \label{eq:snr_link} \\
R_\ell  &= \frac{1}{8}\,\text{MCS\_lookup}(\text{SNR}_\ell,\; \text{standard},\; W)
  \label{eq:rate_link}
\end{align}

This rate $R_\ell$ (in MB/s) is what is assigned to \texttt{link.bandwidth}
before the simulation begins.

\textbf{Explicit bandwidth override:} If a link specifies an explicit
\texttt{bandwidth} value in the scenario YAML, that value is preserved and
the RF-derived rate is not applied.  This allows mixing wired and wireless
links in the same scenario.

%% ============================================================
\section{Carrier Sensing Range}
\label{sec:csrange}

The carrier sensing range $d_\text{CS}$ is the maximum distance at which a
transmission is detected by the Clear Channel Assessment (CCA) mechanism:

\begin{equation}
\label{eq:csrange}
P_\text{tx} - \text{PL}(d_\text{CS}) = \theta_\text{CCA}
\end{equation}

Solving for $d_\text{CS}$:

\begin{equation}
\label{eq:csrange_solve}
d_\text{CS} = d_0 \cdot 10^{\,\frac{P_\text{tx} - \theta_\text{CCA} - \text{PL}(d_0)}{10\,n}}
\end{equation}

With default parameters ($P_\text{tx} = 20\dBm$, $\theta_\text{CCA} =
-82\dBm$, $f = 5\GHz$, $n = 3$), the carrier sensing range is approximately
$71.2$\,m.

Note that the carrier sensing range is determined by the CCA threshold,
while the communication range is determined by the receiver sensitivity
(noise floor plus minimum SNR threshold).  Because the CCA threshold
($-82\dBm$) is higher (less sensitive) than the receiver sensitivity
($-95 + 5 = -90\dBm$), the carrier sensing range can be \emph{shorter}
than the communication range.  This is a well-known property of 802.11
systems and is one source of hidden terminal problems.

%% ============================================================
\section{Conflict Graph}
\label{sec:conflict}

The conflict graph $G_C = (L, E_C)$ is defined over the set of links $L$.
Two links $\ell_a$ and $\ell_b$ are connected by an edge in $E_C$ (i.e.,
they \emph{conflict}) if carrier sensing prevents them from transmitting
simultaneously.  This follows the \emph{protocol interference model}.

\subsection{Conflict Definition}

Let $\text{tx}(\ell)$ and $\text{rx}(\ell)$ denote the transmitter and
receiver nodes of link $\ell$, respectively.  Let
$d(u, v)$ denote the Euclidean distance between nodes $u$ and $v$.

\subsubsection{Without RTS/CTS}

Links $\ell_a$ and $\ell_b$ conflict if either transmitter can sense any node
of the other link:

\begin{equation}
\label{eq:conflict_nortscts}
(\ell_a, \ell_b) \in E_C \iff
\exists\, (u, v) \in
  \bigl\{\text{tx}(\ell_a)\bigr\} \times \bigl\{\text{tx}(\ell_b), \text{rx}(\ell_b)\bigr\}
  \;\cup\;
  \bigl\{\text{tx}(\ell_b)\bigr\} \times \bigl\{\text{tx}(\ell_a), \text{rx}(\ell_a)\bigr\}
\;:\; d(u, v) \leq d_\text{CS}
\end{equation}

In other words, only \emph{transmitters} perform carrier sensing.  The
receiver does not sense the channel before the transmitter initiates.

\subsubsection{With RTS/CTS}

When RTS/CTS is enabled, the virtual carrier sensing mechanism protects both
transmitters and receivers.  Any node of one link sensing any node of another
link creates a conflict:

\begin{equation}
\label{eq:conflict_rtscts}
(\ell_a, \ell_b) \in E_C \iff
\exists\, (u, v) \in
  \bigl\{\text{tx}(\ell_a), \text{rx}(\ell_a)\bigr\} \times
  \bigl\{\text{tx}(\ell_b), \text{rx}(\ell_b)\bigr\}
\;:\; u \neq v \;\wedge\; d(u, v) \leq d_\text{CS}
\end{equation}

RTS/CTS reduces hidden terminal problems by extending the conflict graph, at
the cost of reduced spatial reuse.

\subsection{Maximum Clique Computation}
\label{sec:clique}

For the \texttt{csma\_clique} model (\Cref{sec:clique_model}), each link
$\ell$ needs the size of the largest clique containing it in the conflict
graph:

\begin{equation}
\omega(\ell) = \max\bigl\{|C| : C \text{ is a clique in } G_C \text{ and } \ell \in C\bigr\}
\end{equation}

The implementation uses:
\begin{itemize}[itemsep=2pt]
\item \textbf{Bron-Kerbosch with pivoting} (exact) for networks with $\leq 50$ links.
\item \textbf{Greedy approximation} for larger networks: for each link $\ell$,
  greedily grow a clique from $\ell$ by adding neighbors that are pairwise
  connected.
\end{itemize}

%% ============================================================
\section{Bianchi's CSMA/CA MAC Efficiency Model}
\label{sec:bianchi}

The MAC-layer efficiency of IEEE~802.11 DCF is modeled using Bianchi's
saturation throughput analysis~\cite{bianchi2000}.  This gives the fraction
of channel time used for successful payload delivery when $n$ stations
contend.

\subsection{Coupled Equations}

Bianchi's model solves two coupled nonlinear equations for the transmission
probability $\tau$ and the conditional collision probability $p$:

\begin{align}
\tau &= \frac{2(1 - 2p)}{(1 - 2p)(W_{\min} + 1) + p\,W_{\min}\bigl(1 - (2p)^m\bigr)}
  \label{eq:bianchi_tau} \\
p    &= 1 - (1 - \tau)^{n-1}
  \label{eq:bianchi_p}
\end{align}

where:
\begin{itemize}[itemsep=2pt]
\item $W_{\min} = 16$ is the minimum contention window (CWmin for 802.11),
\item $m = 6$ is the maximum backoff stage ($\text{CW}_\text{max} = W_{\min} \cdot 2^m = 1024$),
\item $n$ is the number of contending stations.
\end{itemize}

These equations are solved by iterative fixed-point iteration, initialized
with $\tau = 2/(W_{\min} + 1)$ and iterated until convergence
($|\tau_\text{new} - \tau| < 10^{-10}$, max 200 iterations).

\subsection{Throughput Efficiency}

From the converged $\tau$, three probabilities characterize each virtual
time slot:

\begin{align}
P_\text{idle}      &= (1 - \tau)^n \label{eq:p_idle} \\
P_\text{success}   &= n\,\tau\,(1 - \tau)^{n-1} \label{eq:p_success} \\
P_\text{collision} &= 1 - P_\text{idle} - P_\text{success} \label{eq:p_collision}
\end{align}

The expected duration of a virtual slot is:

\begin{equation}
\label{eq:eslot}
\mathbb{E}[T_\text{slot}] =
  P_\text{idle} \cdot T_\text{idle}
  + P_\text{success} \cdot T_\text{success}
  + P_\text{collision} \cdot T_\text{collision}
\end{equation}

with the following slot durations (5\,GHz OFDM PHY):

\begin{table}[h]
\centering
\caption{Bianchi model slot duration parameters.}
\label{tab:bianchi_slots}
\begin{tabular}{@{}llr@{}}
\toprule
\textbf{Parameter} & \textbf{Description} & \textbf{Value ($\mu$s)} \\
\midrule
$T_\text{idle}$      & Empty slot duration       & 9 \\
$T_\text{success}$   & Successful frame (DATA + SIFS + ACK) & 500 \\
$T_\text{collision}$ & Collision duration         & 500 \\
\bottomrule
\end{tabular}
\end{table}

The MAC efficiency (fraction of time carrying successful payload) is:

\begin{equation}
\label{eq:eta}
\eta(n) = \frac{P_\text{success} \cdot T_\text{success}}{\mathbb{E}[T_\text{slot}]}
\end{equation}

\textbf{Properties:}
\begin{itemize}[itemsep=2pt]
\item $\eta(1) = 1.0$ (single station, no contention).
\item $\eta(n)$ is monotonically decreasing in $n$.
\item $\eta(n) > 0$ for all $n \geq 1$.
\item Typical values: $\eta(2) \approx 0.88$, $\eta(5) \approx 0.71$,
  $\eta(10) \approx 0.59$.
\end{itemize}

The implementation precomputes $\eta(n)$ for $n = 1, \ldots, 100$ in a
lazy-initialized lookup table.

%% ============================================================
\section{Interference Model Variants}
\label{sec:interference}

The \ncsim{} execution engine computes the effective bandwidth for each link as:

\begin{equation}
\label{eq:engine}
B_\text{eff}(\ell) = \frac{B_\ell \cdot f(\ell)}{N_\ell}
\end{equation}

where $B_\ell$ is the link's base bandwidth (\texttt{link.bandwidth}),
$f(\ell) \in (0, 1]$ is the interference factor, and $N_\ell$ is the number
of concurrent flows sharing link $\ell$.  The two WiFi model variants differ
in how $B_\ell$ and $f(\ell)$ are set.

\subsection{Variant 1: Static Clique Model (\texttt{csma\_clique})}
\label{sec:clique_model}

The static clique model computes all contention effects at setup time and
bakes them into the link bandwidth:

\begin{equation}
\label{eq:clique_bw}
B_\ell = \frac{R_\ell}{\omega(\ell)}
\end{equation}

where $R_\ell$ is the SNR-derived PHY rate (from \Cref{sec:phyrate}) and
$\omega(\ell)$ is the maximum clique size containing $\ell$ (from
\Cref{sec:clique}).

The interference factor is always 1.0:

\begin{equation}
f_\text{clique}(\ell) = 1.0
\end{equation}

\textbf{Rationale:} In the worst case, all links in the maximum clique are
simultaneously active, and CSMA/CA time-sharing gives each an equal fraction.
The clique bound provides a conservative upper bound on the number of mutually
contending links.  This is a static, worst-case analysis independent of which
links are actually active.

\textbf{Properties:}
\begin{itemize}[itemsep=2pt]
\item No dynamic recalculation needed (\texttt{get\_affected\_links} returns $\emptyset$).
\item Conservative: may underestimate throughput when not all clique members are active.
\item Computationally cheap at runtime.
\end{itemize}

\subsection{Variant 2: Dynamic Bianchi Model (\texttt{csma\_bianchi})}
\label{sec:bianchi_model}

The dynamic model correctly separates two distinct interference mechanisms:

\begin{enumerate}
\item \textbf{Contention domain} (conflict graph neighbors): CSMA prevents
  simultaneous transmission.  These links share airtime via Bianchi's model.
\item \textbf{Hidden terminals} (active links \emph{not} in the conflict graph):
  These may transmit simultaneously, causing SINR degradation at the receiver.
\end{enumerate}

\subsubsection{Setup}

The link bandwidth is set to the SNR-only PHY rate (no contention reduction):

\begin{equation}
B_\ell = R_\ell
\end{equation}

\subsubsection{Dynamic Interference Factor}

Let $\mathcal{A}$ be the set of currently active links.  For link $\ell$,
define:

\begin{align}
\mathcal{C}_\ell &= \mathcal{A} \cap \mathcal{N}(\ell)
  && \text{(active contending links---conflict graph neighbors)} \label{eq:contending} \\
\mathcal{H}_\ell &= (\mathcal{A} \setminus \mathcal{N}(\ell)) \setminus \{\ell\}
  && \text{(active hidden terminals---not in conflict graph)} \label{eq:hidden}
\end{align}

where $\mathcal{N}(\ell)$ is the set of conflict graph neighbors of $\ell$.

\paragraph{SINR factor (hidden terminals).}
If $\mathcal{H}_\ell \neq \emptyset$, the SINR at the receiver of $\ell$ is
computed per \Cref{eq:sinr}, with interference from hidden terminal
transmitters only.  The degraded rate is looked up from the MCS table:

\begin{equation}
\label{eq:sinr_factor}
f_\text{SINR}(\ell) = \frac{R_\text{SINR}(\ell)}{R_\ell}
\end{equation}

where $R_\text{SINR}(\ell) = \text{MCS\_lookup}(\text{SINR}_\ell)$ and
$R_\ell$ is the base SNR-only rate.  If no hidden terminals are active,
$f_\text{SINR} = 1.0$.  If the SINR drops below the minimum MCS threshold,
$f_\text{SINR} = 0.01$ (link effectively destroyed).

\paragraph{Contention factor (Bianchi).}
The number of contending stations is $n = 1 + |\mathcal{C}_\ell|$ (the link
itself plus its active contending neighbors).  Each station gets a fraction
$\eta(n)/n$ of the channel:

\begin{equation}
\label{eq:contention_factor}
f_\text{contention}(\ell) = \frac{\eta(n)}{n}
\end{equation}

\paragraph{Combined factor.}

\begin{equation}
\label{eq:combined_factor}
\boxed{f_\text{bianchi}(\ell) = f_\text{SINR}(\ell) \cdot f_\text{contention}(\ell)
= \frac{R_\text{SINR}(\ell)}{R_\ell} \cdot \frac{\eta(n)}{n}}
\end{equation}

clamped to $[0.01, 1.0]$.

\subsubsection{Dynamic Recalculation}

When a transfer starts or completes on a link, all active conflict graph
neighbors are recalculated:

\begin{equation}
\texttt{get\_affected\_links}(\ell) = \mathcal{N}(\ell) \cap \mathcal{A}
\end{equation}

This triggers re-evaluation of the interference factor and rescheduling of
completion events for affected transfers.

\subsection{Comparison of Variants}

\begin{table}[h]
\centering
\caption{Comparison of WiFi interference model variants.}
\label{tab:comparison}
\begin{tabular}{@{}lll@{}}
\toprule
\textbf{Property} & \textbf{\texttt{csma\_clique}} & \textbf{\texttt{csma\_bianchi}} \\
\midrule
Link bandwidth setup & $R_\ell / \omega(\ell)$ & $R_\ell$ \\
Interference factor  & $1.0$ (static) & $\frac{R_\text{SINR}}{R_\ell} \cdot \frac{\eta(n)}{n}$ (dynamic) \\
Contention model     & Worst-case clique bound & Bianchi MAC + active count \\
Hidden terminals     & Not modeled & SINR degradation \\
Recalculation trigger & None & Active neighbor changes \\
Runtime cost         & $O(1)$ per event & $O(|\mathcal{N}(\ell)|)$ per event \\
\bottomrule
\end{tabular}
\end{table}

%% ============================================================
\section{Integration with the Execution Engine}
\label{sec:integration}

The WiFi model integrates with \ncsim{} via two existing abstractions,
requiring \emph{zero changes} to the core execution engine:

\begin{enumerate}
\item \textbf{\texttt{link.bandwidth}} is overwritten at setup time with the
  WiFi-derived PHY rate (or PHY rate / clique size for \texttt{csma\_clique}).

\item \textbf{\texttt{InterferenceModel.get\_interference\_factor()}} returns
  the dynamic factor $f(\ell)$, which the engine multiplies with the link
  bandwidth before applying per-link fair sharing:

\begin{lstlisting}[caption={Engine bandwidth computation (execution\_engine.py).}]
def _get_link_bandwidth(self, link_id, ...):
    link = self.network.links[link_id]
    if self.interference_model is None:
        return link.bandwidth
    active_links = self._get_active_link_ids()
    factor = self.interference_model.get_interference_factor(
        link_id, active_links, self.network
    )
    return link.bandwidth * factor
\end{lstlisting}
\end{enumerate}

The per-flow effective bandwidth is then:
\begin{equation}
B_\text{per\_flow}(\ell) = \frac{B_\ell \cdot f(\ell)}{N_\ell}
\end{equation}

Transfer completion time is:
\begin{equation}
t_\text{complete} = t_\text{now} + \frac{D_\text{remaining}}{B_\text{per\_flow}} + L_\text{total}
\end{equation}

where $D_\text{remaining}$ is the remaining data in MB and $L_\text{total}$
is the total link latency.

%% ============================================================
\section{Worked Example}
\label{sec:example}

Consider the \texttt{wifi\_test.yaml} scenario with default RF parameters:

\begin{itemize}[itemsep=2pt]
\item Four nodes: $n_0(0,0)$, $n_1(30,0)$, $n_2(0,30)$, $n_3(30,30)$.
\item Two links: $\ell_{01}: n_0 \to n_1$ and $\ell_{23}: n_2 \to n_3$.
\item Each link carries a 50\,MB transfer.
\end{itemize}

\subsection{PHY Rate Computation}

For link $\ell_{01}$ with $d = 30$\,m:

\begin{align*}
\text{PL}(d_0 = 1\text{m}) &= 20\log_{10}\!\left(\frac{4\pi \cdot 1 \cdot 5\times 10^9}{3\times 10^8}\right) \approx 46.4\dB \\
\text{PL}(30\text{m}) &= 46.4 + 10 \cdot 3 \cdot \log_{10}(30) \approx 46.4 + 44.3 = 90.7\dB \\
P_\text{rx} &= 20 - 90.7 = -70.7\dBm \\
\text{SNR} &= -70.7 - (-95) = 24.3\dB \\
\text{MCS lookup} &\Rightarrow \text{MCS 5 (64-QAM 2/3), min SNR = 22\dB}: 68.8\Mbps \\
R_{\ell_{01}} &= 68.8 / 8 = 8.6\MBps
\end{align*}

Both links have identical geometry, so $R_{\ell_{23}} = 8.6\MBps$.

\subsection{Conflict Graph}

Carrier sensing range: $d_\text{CS} = 71.2$\,m.  The transmitter $n_0$ is at
distance $30$\,m from $n_2$ and $42.4$\,m from $n_3$; both within $d_\text{CS}$.
Therefore $\ell_{01}$ and $\ell_{23}$ conflict.  $\omega(\ell_{01}) =
\omega(\ell_{23}) = 2$.

\subsection{\texttt{csma\_clique} Result}

\begin{align*}
B_{\ell_{01}} &= 8.6 / 2 = 4.3\MBps, \quad f = 1.0 \\
t_\text{transfer} &= 50 / 4.3 = 11.63\text{s} \\
\text{Makespan} &= 0.01\text{s (task)} + 11.63\text{s (transfer)} + 0.01\text{s (task)} = 11.65\text{s}
\end{align*}

\subsection{\texttt{csma\_bianchi} Result}

\begin{align*}
B_{\ell_{01}} &= 8.6\MBps \text{ (base rate)} \\
n &= 2 \text{ contending stations} \\
\eta(2) &\approx 0.88, \quad f_\text{contention} = 0.88 / 2 = 0.44 \\
f_\text{SINR} &= 1.0 \text{ (no hidden terminals---both links are in conflict graph)} \\
f &= 1.0 \times 0.44 = 0.44 \\
B_\text{eff} &= 8.6 \times 0.44 = 3.78\MBps \\
t_\text{transfer} &= 50 / 3.78 = 13.21\text{s} \\
\text{Makespan} &\approx 13.22\text{s}
\end{align*}

The Bianchi model gives a longer transfer time because $\eta(2)/2 \approx
0.44 < 1/2 = 0.50$ (the clique model's share).  The difference accounts for
MAC protocol overhead (backoff slots, collisions) that the clique model's
simple $1/\omega$ bound does not capture.

%% ============================================================
\section{YAML Configuration}
\label{sec:yaml}

\begin{lstlisting}[language={},caption={WiFi scenario YAML example.}]
scenario:
  name: "WiFi CSMA Bianchi Test"
  network:
    nodes:
      - {id: n0, compute_capacity: 1000, position: {x: 0, y: 0}}
      - {id: n1, compute_capacity: 1000, position: {x: 30, y: 0}}
    links:
      # No bandwidth -- derived from RF model
      - {id: l01, from: n0, to: n1, latency: 0.0}
      # Explicit bandwidth preserved as override
      - {id: l23, from: n2, to: n3, bandwidth: 50, latency: 0.0}
  config:
    interference: csma_bianchi   # or csma_clique
    rf:
      tx_power_dBm: 20
      freq_ghz: 5.0
      path_loss_exponent: 3.0
      noise_floor_dBm: -95
      cca_threshold_dBm: -82
      channel_width_mhz: 20
      wifi_standard: "ax"
      shadow_fading_sigma: 0.0
      rts_cts: false
\end{lstlisting}

%% ============================================================
\section{Summary}
\label{sec:summary}

The 802.11 WiFi model in \ncsim{} provides a physically-grounded replacement
for simple heuristic interference models.  The key contributions are:

\begin{enumerate}
\item \textbf{Distance-derived link rates} from log-distance path loss and
  802.11n/ac/ax MCS tables, eliminating the need to manually specify wireless
  link bandwidths.
\item \textbf{Protocol-model conflict graph} from carrier sensing range, with
  optional RTS/CTS for receiver protection.
\item \textbf{Two complementary interference variants}: a fast static clique
  bound and a dynamic Bianchi-based model that captures both CSMA contention
  and hidden terminal SINR degradation.
\item \textbf{Zero engine modifications}: the model composes cleanly with
  \ncsim{}'s existing \texttt{link.bandwidth} and \texttt{InterferenceModel}
  abstractions.
\end{enumerate}

\begin{thebibliography}{9}

\bibitem{bianchi2000}
G.~Bianchi,
``Performance Analysis of the IEEE 802.11 Distributed Coordination Function,''
\emph{IEEE Journal on Selected Areas in Communications},
vol.~18, no.~3, pp.~535--547, March 2000.

\bibitem{gupta2000}
P.~Gupta and P.~R.~Kumar,
``The Capacity of Wireless Networks,''
\emph{IEEE Transactions on Information Theory},
vol.~46, no.~2, pp.~388--404, March 2000.

\bibitem{garetto2008}
M.~Garetto, T.~Salonidis, and E.~W.~Knightly,
``Modeling Per-Flow Throughput and Capturing Starvation in CSMA Multi-Hop
Wireless Networks,''
\emph{IEEE/ACM Transactions on Networking},
vol.~16, no.~4, pp.~864--877, August 2008.

\bibitem{ieee80211ax}
IEEE Std 802.11ax-2021,
``IEEE Standard for Information Technology---Telecommunications and
Information Exchange Between Systems---Local and Metropolitan Area
Networks---Specific Requirements---Part 11: Wireless LAN Medium Access Control
(MAC) and Physical Layer (PHY) Specifications---Amendment 1: Enhancements for
High-Efficiency WLAN,'' 2021.

\end{thebibliography}

\end{document}
